\chapter{Architektura i struktura aplikacji}
\label{cha:architektura}

\section{Architektura aplikacji}
\label{sec:architektura}

Aplikacja ma architekturę warstwową. Wyróżniono cztery warstwy o odrębnych
odpowiedzialnościach:

\begin{enumerate}
  \item \textbf{Warstwa modelu (\texttt{Models})} - klasy encji odwzorowujące obiekty dziedziny
        (użytkownik, grupa, test, pytanie, podejście, odpowiedź), mapowane na tabele bazy danych
        przez Entity Framework Core.
  \item \textbf{Warstwa dostępu do danych (\texttt{Data})} - kontekst bazy danych
        (\texttt{AppDbContext}) z konfiguracją relacji i dziedziczenia, mechanizm migracji oraz klasa
        zasilająca bazę danymi początkowymi (\texttt{DbSeeder}).
  \item \textbf{Warstwa logiki i usług (\texttt{Services})} - repozytorium operacji bazodanowych
        (\texttt{DbRepository}), serwis oceniania (\texttt{GradingService}), eksport do CSV
        (\texttt{CsvExporter}), haszowanie haseł (\texttt{PasswordHash}) oraz wspólna obsługa
        wyjątków (\texttt{Safe}). Tu znajdują się także obiekty transferu danych (DTO).
  \item \textbf{Warstwa prezentacji (\texttt{Views})} - okna i ekrany aplikacji zbudowane w
        technologii WPF, odpowiedzialne za wyświetlanie danych i obsługę interakcji użytkownika.
\end{enumerate}

Przepływ jest jednokierunkowy: widok wywołuje warstwę usług, ta korzysta z repozytorium, a
repozytorium komunikuje się z bazą przez kontekst Entity Framework Core. Widoki nigdy nie odwołują
się do bazy bezpośrednio - zawsze pośredniczy w tym warstwa usług.

\section{Struktura folderów i organizacja kodu}
\label{sec:foldery}

Kod źródłowy podzielono na foldery odpowiadające warstwom aplikacji (listing~\ref{lst:foldery}).

\begin{lstlisting}[caption={Struktura katalogów projektu.}, label=lst:foldery]
ProjektPO2/
|-- App.xaml(.cs)              // punkt startowy aplikacji
|-- MainWindow.xaml(.cs)       // powloka: panel boczny + obszar tresci
|-- Theme.xaml                 // motyw: kolory, style kontrolek
|-- Models/                    // encje (User, Quiz, Question, ...)
|-- Data/                      // AppDbContext, DbSeeder, fabryka kontekstu
|-- Migrations/                // migracje EF Core
|-- Services/                  // DbRepository, GradingService, CsvExporter,
|                              //   PasswordHash, Safe, DTO (AppModels)
|-- Views/                     // ekrany aplikacji (Login, Teacher*, Quiz, ...)
\-- Commands/                  // RelayCommand
\end{lstlisting}

\section{Zastosowane wzorce projektowe}
\label{sec:wzorce}

W projekcie wykorzystano kilka rozwiązań i wzorców porządkujących kod:

\begin{itemize}
  \item \textbf{Architektura warstwowa} - oddzielenie modelu, dostępu do danych, logiki i prezentacji.
  \item \textbf{Repozytorium (Repository)} - klasa \texttt{DbRepository} hermetyzuje cały dostęp do
        bazy danych i odcina od niej resztę aplikacji.
  \item \textbf{Obiekty transferu danych (DTO)} - lekkie klasy w warstwie usług oddzielają dane
        prezentowane w widokach od encji bazodanowych.
  \item \textbf{Dziedziczenie i polimorfizm} - trzy hierarchie klas abstrakcyjnych (użytkownicy,
        pytania, odpowiedzi) mapowane na bazę strategią Table-per-Hierarchy.
  \item \textbf{Interfejs \texttt{IPublishable}} - implementowany przez klasę \texttt{Quiz},
        ujednolica sterowanie widocznością testu (metody \texttt{Publish()} i \texttt{Unpublish()}).
\end{itemize}

\section{Mechanizmy walidacji danych}
\label{sec:walidacja}

Walidacja realizowana jest na trzech poziomach:

\begin{enumerate}
  \item \textbf{Interfejs użytkownika} - pola formularzy są sprawdzane przed zapisem: niepuste
        wymagane pola, liczby ograniczane do dozwolonego zakresu (np. czas, próg, liczba podejść),
        a edytor testu blokuje zapis testu niekompletnego (pytanie bez poprawnej odpowiedzi, puste
        pola, brak przypisanej grupy).
  \item \textbf{Logika aplikacji} - serwis oceniania pilnuje zakresu przyznawanych punktów, a
        repozytorium dba o spójność operacji (np. bezpieczną kolejność usuwania powiązanych danych).
  \item \textbf{Baza danych} - ograniczenia na poziomie tabel (klucze obce, unikalny indeks)
        zapewniają spójność nawet przy bezpośrednim dostępie do bazy.
\end{enumerate}

\section{Obsługa błędów i wyjątków}
\label{sec:bledy}

Operacje na bazie danych oraz operacje wejścia/wyjścia są wykonywane za pośrednictwem wspólnej klasy
\texttt{Safe}, która przechwytuje wyjątki i prezentuje użytkownikowi czytelny komunikat zamiast
powodować awarię aplikacji. Dzięki temu np. brak połączenia z bazą czy błąd zapisu pliku CSV nie
zatrzymuje programu, a użytkownik dowiaduje się o przyczynie problemu.
